%=============================================================================!
% 2.1  FORCE AND THE FIRST LAW                                                !
%=============================================================================!
\section{Force and the first law}
\label{sec:ne-force}

In everyday speech a force is a push or a pull: a hand pushing a door, a
rope pulling a sledge, a magnet pulling a pin. To build a law of motion on
forces we first need to measure them, and this section shows how, before
stating the first of Newton's three laws.

\subsection*{Measuring a force}

A spring stretches when it is pulled, and the harder the pull, the further
it stretches. A spring with a scale beside it, called a spring balance,
therefore turns a pull into a reading. We can define a unit of force as
the pull that stretches one particular standard spring by one particular
amount. Two identical standard springs pulling side by side, each
stretched by that amount, then exert two units, three springs exert three
units, and so on (\cref{fig:ne-force}a). In this way a force acquires a
size, a number of units. \Cref{sec:ne-second} replaces this home-made
unit by the
standard one, the newton.

A force also has a direction, the direction in which it pulls, so it is a
vector. When several forces act on one body together, experiment shows
that their combined effect on its motion from place to place is that of a
single force equal to their vector sum, called the net force,
\begin{equation*}
   \vb{F} = \vb{F}_1 + \vb{F}_2 + \dotsb ,
\end{equation*}
in the sense in which a net profit is what remains when everything is
added up. A pull of 4 units towards the east and a pull of 3 units towards
the north, for example, are the vectors $(4, 0)$ and $(0, 3)$ when the
$x$ axis points east and the $y$ axis north. Adding component by
component (\cref{sec:mo-plane}) gives $(4, 3)$. This is the same as
drawing the second arrow from the tip of the first, since starting at the
tip of $(4, 0)$ and moving by $(0, 3)$ ends at $(4, 3)$
(\cref{fig:ne-force}b). The two pulls therefore act like a single pull of
length $(4^2 + 3^2)^{1/2} = 5$ units, by Pythagoras' theorem, directed at
an angle $\theta$ north of east. Writing $(4, 3) = 5(\cos\theta,
\sin\theta)$ gives $\tan\theta = \sin\theta/\cos\theta = 3/4$. The tangent can be
undone as the sine was in \cref{sec:mo-trig}: for any number $s$,
$\arctan s$ is the angle between $-\pi/2$ and $\pi/2$ whose tangent is
$s$, and here $\theta = \arctan(3/4) = \SI{0.644}{\radian}$, or
\ang{36.9}. Two forces of equal size pulling
in opposite directions add to zero, so a body can be pushed and pulled by
several forces while the net force on it vanishes. Such a body may
still turn: two equal and opposite pushes at opposite corners of a book
lying on a table add to zero, yet set the book spinning, because where a
force acts decides how it turns a body. Turning is the subject of
Chapter~5, and until then every body is a particle, for which only the
net force matters.

\begin{figure}[htb]
\centering
\begin{tikzpicture}[line cap=round, line join=round,
   force/.style={-{Stealth[length=5pt]}, line width=1.1pt},
   coil/.style={decorate, decoration={coil, aspect=0.45,
      segment length=3pt, amplitude=3pt, pre length=3pt,
      post length=3pt}, line width=0.6pt}]
   % (a) one spring, then two side by side, each stretched by one unit
   \begin{scope}
      \node[anchor=west] at (-0.4,1.75) {\small (a)};
      \fill[black!35] (-0.25,1.45) rectangle (2.95,1.55);
      \draw[coil] (0.3,1.45) -- (0.3,0.35);
      \draw[black!40] (0.0,0.35) -- (0.0,1.25);
      \foreach \y in {0.35,0.65,0.95,1.25}
         \draw[black!40] (0.0,\y) -- (0.08,\y);
      \draw[force] (0.3,0.35) -- (0.3,-0.35);
      \node[below] at (0.3,-0.38) {\small 1 unit};
      \draw[coil] (1.9,1.45) -- (1.9,0.35);
      \draw[coil] (2.4,1.45) -- (2.4,0.35);
      \draw[line width=0.8pt] (1.8,0.35) -- (2.5,0.35);
      \draw[force] (2.15,0.35) -- (2.15,-1.05);
      \node[below] at (2.15,-1.08) {\small 2 units};
   \end{scope}
   % (b) two forces added tip to tail
   \begin{scope}[xshift=5.0cm, yshift=-0.9cm, scale=0.62]
      \node[anchor=west] at (-0.6,4.1) {\small (b)};
      \draw[force] (0,0) -- (4,0) node[midway, below] {\small 4 east};
      \draw[force] (4,0) -- (4,3) node[midway, right] {\small 3 north};
      \draw[force, accent] (0,0) -- (4,3)
         node[midway, above left] {\small net force, 5};
      \draw[black!50, line width=0.4pt] (1.0,0) arc[start angle=0,
         end angle=36.87, radius=1.0];
      \node at (1.45,0.42) {\small \ang{36.9}};
      \fill (0,0) circle (1.5pt);
   \end{scope}
\end{tikzpicture}
\caption{Measuring and adding forces. (a) A unit of force stretches a
standard spring by a set amount; two such springs pulling together exert
two units. (b) Pulls of 4 units east and 3 units north, drawn tip to tail,
add to a net force of 5 units at \ang{36.9} north of east.}
\label{fig:ne-force}
\end{figure}

\subsection*{The first law}

A book slid across a table slows down and stops. The table and the air
exert forces on it: friction, the grip of one surface sliding over
another, and drag, the resistance of the air or liquid a body moves
through. Where these forces are made small, the book's behaviour changes.
On smooth ice, or on an air-hockey table, where a cushion of air lifts the
puck clear of the surface, a puck glides on at nearly constant speed in a
straight line. We call a surface smooth when the friction it exerts is
small enough to neglect.

Since forces acting together behave like their net force, a body whose
forces add to zero moves as one with no force on it at all. Newton's
first law states the result: a body on which no net force acts keeps moving in a straight line at constant
speed, that is, at constant velocity, and a body at rest stays at rest. A
force is therefore not needed to keep a body moving, only to change its
motion. The next section shows that this law cannot be true for every
observer, and uses that fact to pick out the viewpoints from which all
three laws hold.

\begin{selfcheck}
A puck glides across smooth ice in a straight line at a steady
\SI{3}{\metre\per\second}. What is the net force on it?
\end{selfcheck}
